%!TEX root =../mapp-challenge-18-game-book.tex
% ^ leave for LaTeXTools build functionality

\phChapterWorksheet{Main Puzzle ?}{Doppl's Unending Enigma}

\newcommand{\mappDoppl}{\includegraphics[height=0.12in]{assets/doppl.pdf}}
\newcommand{\mappDopplOmega}{\includegraphics[height=0.12in]{assets/doppl-omega.pdf}}
\newcommand{\mappDopplOmegaSquared}{\includegraphics[height=0.12in]{assets/doppl-omega2.pdf}}

It's time to put your training to the test by trying to befriend \textbf{Doppl}!
% This Ordinary-Type \mappMobimon{} multiplies very quickly, so it's
% considered to be quite hard to Wrangle.
As you might expect, the natural
numbers, known as the \textbf{finite ordinals} to \mappMobimon{} Wranglers,
are used to account for a finite amount of Doppl.

\begin{itemize}
\item \(0=\) no Doppl

\item \(1=\) \mappDoppl

\item \(2=\) \mappDoppl\mappDoppl

\item \(3=\) \mappDoppl\mappDoppl\mappDoppl

\item \(4=\) \mappDoppl\mappDoppl\mappDoppl\mappDoppl
\end{itemize}

However, the problem is that \textbf{infinite} groups of Doppl are quite
common out in the wild! Wranglers use the lowercase Greek letter \(\omega\)
("omega") to represent the shortest possible chain of infinitely-many Doppl.

\begin{itemize}
  \item \(\omega=\)
  \mappDoppl\mappDoppl\mappDoppl\mappDoppl\mappDoppl\mappDoppl\mappDoppl
  \mappDoppl\mappDoppl\mappDoppl\mappDoppl\mappDoppl\mappDoppl\mappDoppl
  \mappDoppl\mappDoppl\mappDoppl\mappDoppl\mappDoppl\mappDoppl\mappDoppl
  \mappDoppl\mappDoppl\mappDoppl\mappDoppl\mappDoppl\mappDoppl\mappDoppl
  \mappDoppl\mappDoppl\mappDoppl\mappDoppl\mappDoppl\mappDoppl\mappDoppl
  \mappDoppl\mappDoppl\mappDoppl\mappDoppl\mappDoppl\mappDoppl\mappDoppl
  \mappDoppl\mappDoppl\mappDoppl\mappDoppl\mappDoppl\mappDoppl\mappDoppl
\end{itemize}

Of course, if they were all the same size, they would quickly run out
of space for their \mappMobimon{} battles. That's
why groups of Doppl will shrink down as necessary to fit into whatever space
is available. Here's a more typical representation of \(\omega\)-many Doppl.

\begin{itemize}
\item \(\omega=\) \mappDopplOmega
\end{itemize}

This \(\omega\) is the first \textbf{infinite ordinal}, but it's not the last!
You see, bizzare as it sounds, there's always room for another
Doppl to join the party.

\begin{itemize}
\item\(\omega+1=\) \mappDopplOmega\mappDoppl

\item\(\omega+2=\) \mappDopplOmega\mappDoppl\mappDoppl

\item\(\omega+3=\) \mappDopplOmega\mappDoppl\mappDoppl\mappDoppl
\end{itemize}

So \(\omega+1,\omega+2,\omega+3\) are the next three infinte ordinals.
And yes, Wrangers have reported Doppl groups like these as well.

\begin{itemize}
\item \(\omega+\omega=\omega\cdot2=\) \mappDopplOmega\mappDopplOmega

\item \(\omega+\omega+\omega+5=\omega\cdot 3 + 5=\) \mappDopplOmega\mappDopplOmega\mappDopplOmega\mappDoppl\mappDoppl\mappDoppl\mappDoppl\mappDoppl

\item \(\omega+\omega+\omega+\cdots=\omega\cdot\omega=\omega^2=\)
\mappDopplOmega\mappDopplOmega\mappDopplOmega\mappDopplOmega\mappDopplOmega
\mappDopplOmega\mappDopplOmega\mappDopplOmega\mappDopplOmega\mappDopplOmega
\mappDopplOmega\mappDopplOmega\mappDopplOmega\mappDopplOmega\mappDopplOmega
\mappDopplOmega\mappDopplOmega\mappDopplOmega\mappDopplOmega\mappDopplOmega
\mappDopplOmega\mappDopplOmega\mappDopplOmega\mappDopplOmega\mappDopplOmega
\mappDopplOmega\mappDopplOmega\mappDopplOmega\mappDopplOmega\mappDopplOmega
\mappDopplOmega\mappDopplOmega\mappDopplOmega\mappDopplOmega\mappDopplOmega
\mappDopplOmega\mappDopplOmega\mappDopplOmega\mappDopplOmega\mappDopplOmega

= \mappDopplOmegaSquared
\end{itemize}

\newpage

After the discovery of Doppl, the \mappMobimon{} community noticed that
the rules of \textbf{ordinal arithmetic} behave differently when infinite
ordinals are involved. Notice what happens when groups of Doppl are added
together: a finite group of Doppl will attach to an infinite group to its
right.

\begin{itemize}
\item \(3+4=\) \{\mappDoppl\mappDoppl\mappDoppl\}%
\{\mappDoppl\mappDoppl\mappDoppl\mappDoppl\} =
\mappDoppl\mappDoppl\mappDoppl\mappDoppl\mappDoppl\mappDoppl\mappDoppl{}
\(=7\)

\item \(7+\omega=\)
\{\mappDoppl\mappDoppl\mappDoppl\mappDoppl\mappDoppl\mappDoppl\mappDoppl\}%
\{\mappDopplOmega\} =
\mappDoppl\mappDoppl\mappDoppl\mappDoppl\mappDoppl\mappDoppl\mappDoppl%
\mappDopplOmega =
\mappDopplOmega{} =
\(\omega\)

\item \((\omega\cdot4+3)+(\omega\cdot2+5)=\)
\{\mappDopplOmega\mappDopplOmega\mappDopplOmega\mappDopplOmega\mappDoppl
\mappDoppl\mappDoppl\}%
\{\mappDopplOmega\mappDopplOmega\mappDoppl\mappDoppl\mappDoppl\mappDoppl
\mappDoppl\}

\hspace{1.55in} \(=\)
\{\mappDopplOmega\}\{\mappDopplOmega\}\{\mappDopplOmega\}\{\mappDopplOmega\}%
\{\mappDoppl\mappDoppl\mappDoppl\mappDopplOmega\}%
\{\mappDopplOmega\}\mappDoppl\mappDoppl\mappDoppl\mappDoppl\mappDoppl{}

\hspace{1.55in} \(=\omega\cdot6+5\)
\end{itemize}

Doppl are also known to multiply. When this occurs, each Doppl in the second
factor splits into a copy of the Doppl group given by the first factor.
But when infinite groups are around, the normal distribution rules don't apply,
so be careful!

\begin{itemize}
\item \(3\cdot 4=\)
\{\mappDoppl\mappDoppl\mappDoppl\}\{\mappDoppl\mappDoppl\mappDoppl\}\{\mappDoppl\mappDoppl\mappDoppl\}\{\mappDoppl\mappDoppl\mappDoppl\} \(=12\)

\item \((\omega+1)\cdot 2=\)
\{\mappDopplOmega\mappDoppl\}\{\mappDopplOmega\mappDoppl\} =
\{\mappDopplOmega\}\{\mappDoppl\mappDopplOmega\}\mappDoppl{} \(=\omega\cdot 2+1\)

\item \(2\cdot(\omega+1)=\)
\{\mappDoppl\mappDoppl\}%
\resizebox{0.4in}{!}{\{\mappDoppl\mappDoppl\}}%
\resizebox{0.3in}{!}{\{\mappDoppl\mappDoppl\}}%
\resizebox{0.2in}{!}{\{\mappDoppl\mappDoppl\}}%
\resizebox{0.1in}{!}{\{\mappDoppl\mappDoppl\}}%
\{\mappDoppl\mappDoppl\} =
\mappDopplOmega\mappDoppl\mappDoppl \(=\omega+2\)

\item \((\omega+1)\cdot \omega=\)
\{\mappDopplOmega\mappDoppl\}%
\resizebox{0.6in}{!}{\{\mappDopplOmega\mappDoppl\}}%
\resizebox{0.5in}{!}{\{\mappDopplOmega\mappDoppl\}}%
\resizebox{0.4in}{!}{\{\mappDopplOmega\mappDoppl\}}%
\resizebox{0.3in}{!}{\{\mappDopplOmega\mappDoppl\}}%
\resizebox{0.2in}{!}{\{\mappDopplOmega\mappDoppl\}}%
\resizebox{0.1in}{!}{\{\mappDopplOmega\mappDoppl\}}%

\hspace{0.7in}
\(=\)
\{\mappDopplOmega\}\{\mappDoppl\mappDopplOmega\}%
\{\mappDoppl\mappDopplOmega\}\{\mappDoppl\mappDopplOmega\}%
\(\cdots =\)
\mappDopplOmegaSquared \(=\omega^2\)
\end{itemize}

% Finally, \textbf{oridnal exponentiation} follows the formula
% \(\omega^{n+1}=\omega^n\cdot\omega\). That means \(\omega^2\) is formed by
% each Doppl in \(\omega\) splitting into its own \(\omega\),
% and \(\omega^3\) is formed
% by each Doppl in \(\omega\) splitting into \(\omega^2\)...

Can you find the missing values in these Doppl-inspired arithmetic problems?

\begin{align*}
  (\omega+3)\cdot(\omega+5)&=\omega^2\cdot E+\omega\cdot G+S \\
        %M=1,G=5,S=3
  \omega+1+\omega+3+\omega+5+\omega+7&=\omega\cdot I+O \\
        %I=4,O=7
  3\cdot\omega+\omega^2\cdot 5+4\cdot(\omega^2+2)&=
        \omega^2\cdot L+\omega\cdot R+N \\
        %L=6,R=0,N=8
  2\cdot(2+\omega\cdot3)+(\omega\cdot3+2)\cdot2&=
        \omega\cdot A+M
        %A=9,M=2
\end{align*}

If so, you should be able to decode \(24007042951\) into a Doppl's
favorite \mappMobimon{} attack! %solution: MIRRORIMAGE

% Include below for aucTeX integration
%%% Local Variables:
%%% mode: latex
%%% TeX-master: "../mapp-challenge-18-game-book"
%%% End:
